\documentclass[../textbook.tex]{subfiles}
\begin{document}
\bookchapter{信息检索}{ch:information-retrieval}

检索系统需要在规模庞大且不断变化的语料中，找到与当前信息需求有关、可访问并且版本有效的材料。词语是否相同、向量是否接近、排名是否靠前，分别只是相关性的不同近似。系统必须同时回答两个问题：怎样表示并排序候选，以及怎样在有限时间和内存内找到它们。

\term{信息检索}{Information Retrieval}{IR}输出的是有序材料集合。本章讨论词法评分、表示学习、近似索引与检索评估；下游如何将材料组织为生成上下文、如何核对回答与引用，属于检索增强生成的架构问题。高召回可以提供必要材料，却不自动证明答案正确。

\section{检索对象及语料结构}

\subsection{文档、片段及相关性}
令查询为 $q$，检索对象集合为 $\mathcal D$，相关性等级为 $r(q,d)$。这里的对象 $d$ 可以是整篇文档、段落或表格片段，必须与标注和评估一致。一个系统用片段排序、另一个系统按文档统计命中时，直接比较数字可能掩盖重复片段占据前列的问题。

\term{语义切分}{Semantic Chunking}{}依据标题、段落、句子、代码或表格结构形成边界，再结合长度限制组织片段。较短片段有利于定位，但可能丢失否定条件、表头和适用范围；较长片段提供上下文，却可能稀释主题并增加编码成本。重叠窗口减少边界遗漏，也增加重复和索引体积，因此应同时计量召回与重复率。

解析阶段应保留标题层级、页码或源偏移、版本、生效时间、语言与访问条件。表格单元脱离表头常失去意义，扫描文档的文字识别错误又会影响精确词匹配。解析与文字识别质量决定后续检索可用的信息。原始文档、派生片段与索引条目之间的血缘关系，遵循\bookxref{ch:data-engineering}所述原则。

\begin{table}[htbp]
\centering
\caption{主要符号与形状。}
\label{tab:information-retrieval-symbols}
\begin{tabularx}{\linewidth}{@{}lX@{}}
\toprule
符号 & 含义\\
\midrule
$N,n_t,f(t,d)$ & 检索对象总数、包含词项 $t$ 的对象数、对象内词频\\
$|d|,\bar\ell$ & 对象的分析器词项长度与平均长度\\
$k_1,b$ & BM25 饱和参数与长度归一化参数\\
$u,v_d\in\Real^h$ & 查询与文档的 $h$ 维表示\\
$\tau,s(q,d)$ & 对比学习温度与检索分数\\
$K,R_q^{(K)},G_q$ & 截断长度、前 $K$ 个结果、相关对象集合\\
$M,m,c$ & 图邻接预算、PQ 子空间数、RRF 平滑常数\\
\bottomrule
\end{tabularx}
\end{table}

\subsection{词法分析及匹配对象}
\term{词项}{Term}{}是检索分析器生成的匹配单位，不一定等于语言模型词元。中文分词、大小写、字符规范化、数字编号、同义词与停用词规则都会改变词项集合。查询与文档应使用兼容的分析规则；变更规则一般需要重新构建相应索引。

过度规范化可能破坏型号、错误码和带连字符的标识。对精确编号和普通自然语言，可以使用不同字段及匹配策略，而非强迫所有内容共享同一种切分。字段权重也应作为评分配置，不能在解释分数时忽略。

\section{倒排索引及概率排序}

\subsection{词项候选定位}
\term{倒排索引}{Inverted Index}{}为每个词项保存包含它的对象列表，即\term{倒排列表}{Posting List}{}。每个条目可含文档标识、词频、位置和字段信息。查询不再扫描全部文本，而是读取命中词项的列表，合并候选并计算分数。词项位置可支持短语与邻近关系，只有文档标识和词频则不能直接恢复这些约束。

设查询词项集合为 $T_q$，若完整遍历相关列表，处理量与 $\sum_{t\in T_q}|P_t|$ 有关；常见高频词会产生长列表。跳跃结构和基于分数上界的剪枝可以减少计算，但必须使用有效上界，才能保持指定评分下的精确前 $K$ 名。压缩列表减少内存与读取量，也增加解码代价。

\subsection{概率动机及独立性近似}
概率排序思想是在给定查询和固定损失条件下，优先返回相关概率更高的对象。为说明稀有词为何重要，考虑二元相关变量 $R$ 和词项出现指标 $x_t\in\{0,1\}$。假设给定 $R$ 后词项条件独立，记 $p_t=P(x_t=1\mid R=1)$、$u_t=P(x_t=1\mid R=0)$，则贝叶斯对数比为
\begin{align}
\log\frac{P(R=1\mid x)}{P(R=0\mid x)}
={}&\text{常数}+\sum_t\left[x_t\log\frac{p_t}{u_t}+(1-x_t)\log\frac{1-p_t}{1-u_t}\right]\\
={}&\text{常数}+\sum_tx_t\log\frac{p_t(1-u_t)}{u_t(1-p_t)}.
\end{align}
忽略与候选无关的常数后，词项提供可加的对数证据。缺少相关性标注时，对 $p_t,u_t$ 作近似并使用文档频率估计，会得到与逆文档频率相关的权重。该模型以词项条件独立为近似；BM25 进一步引入词频饱和与文档长度校正\cite{robertson2009bm25}。

\subsection{BM25 的饱和及长度归一化}
\term{最佳匹配评分}{Best Matching 25}{BM25}的一种常用形式为
\begin{equation}
s(q,d)=\sum_{t\in T_q}\operatorname{IDF}(t)
\frac{f(t,d)(k_1+1)}{f(t,d)+k_1(1-b+\frac{b|d|}{\bar\ell})},
\label{eq:ch32-bm25}
\end{equation}
其中\term{逆文档频率}{Inverse Document Frequency}{IDF}取非负变体
\begin{equation}
\operatorname{IDF}(t)=\log\left(1+\frac{N-n_t+0.5}{n_t+0.5}\right).
\end{equation}
本章用自然对数，并将查询重复词去重；其他实现可能保留查询词频因子、采用允许负值的 IDF 或按字段计算长度。公式版本必须固定，不能把不同实现的分数直接比较。

固定长度项 $c_d=k_1(1-b+\frac{b|d|}{\bar\ell})>0$，令 $g(f)=\frac{(k_1+1)f}{f+c_d}$，有
\begin{equation}
g'(f)=\frac{(k_1+1)c_d}{(f+c_d)^2}>0,\qquad
g''(f)=-\frac{2(k_1+1)c_d}{(f+c_d)^3}<0.
\end{equation}
词频贡献递增但增益递减，且 $f\to\infty$ 时趋于 $k_1+1$。在固定 $f>0,b>0$ 时，长度增大提高分母，使分数下降；$b=0$ 不作长度归一化，$b=1$ 使用完整长度比例。长度归一化校正篇幅带来的词项出现机会差异。

\begin{assumption}[BM25 的数值解释]
统计对象与评分对象一致，$N>0,\bar\ell>0,k_1>0,0\le b\le1$。长度使用同一分析器词项数，语料统计绑定一个索引快照。空对象不参与平均长度的异常除法，查询词项规则已固定。
\end{assumption}

\subsection{双文档评分}
\begin{example}
\label{q:ch32:example:01}


设 $N=100$，词项甲出现在 $10$ 个对象，乙出现在 $50$ 个对象。则 $\operatorname{IDF}(\text{甲})=\log(1+\frac{90.5}{10.5})\approx2.2637$，$\operatorname{IDF}(\text{乙})=\log2\approx0.6931$。取 $k_1=1.2,b=0.75,\bar\ell=100$。

文档 A 长 $100$，甲、乙词频为 $(2,1)$，长度项为 $1.2$，两个词频因子为 $\frac{4.4}{3.2}=1.375$ 和 $\frac{2.2}{2.2}=1$，总分约为 $2.2637\times1.375+0.6931=3.8057$。文档 B 长 $200$，词频为 $(3,2)$，长度项为 $2.1$，因子为 $\frac{6.6}{5.1}\approx1.2941$ 和 $\frac{4.4}{4.1}\approx1.0732$，总分约为 $3.6735$。B 词频更高，却因更长而排名稍低。

该分数用于比较候选的词项相关性。语料增量会改变 $N,n_t,\bar\ell$；同一个固定阈值在不同快照或不同字段上可能失去原意义。阈值应在目标任务的开发集上校准，并与无答案情况一起评估。
\end{example}

\section{双编码器及对比学习}

\subsection{独立编码实现离线索引}
\term{稠密检索}{Dense Retrieval}{}使用低维连续表示进行匹配。\term{双编码器}{Dual Encoder}{}分别计算查询向量 $u=f_\theta(q)\in\Real^h$ 和文档向量 $v_d=g_\phi(d)\in\Real^h$，以 $s(q,d)=u^{\mathsf T}v_d$ 等方式评分。文档向量可以离线生成，在线仅编码查询并搜索；稠密段落检索是这一设计的典型研究实例\cite{karpukhin2020dpr}。

查询与文档编码器可以共享或不共享参数，还可能使用不同任务前缀。查询与文档向量需使用相容的编码空间，编码器版本与训练协议共同确定该空间。归一化、池化、截断和编码模板均属于表示协议。

\subsection{对比损失及梯度方向}
给定查询 $q_i$、正例 $d_i^+$ 和候选文档集合 $\mathcal C_i$，\term{对比学习}{Contrastive Learning}{}可采用
\begin{equation}
\mathcal L_i=-\log\frac{\exp(\frac{s(q_i,d_i^+)}{\tau})}{\sum_{d\in\mathcal C_i}\exp(\frac{s(q_i,d)}{\tau})},\qquad\tau>0.
\label{eq:ch32-contrast}
\end{equation}
令候选概率为 $p_{ij}$、正例指示为 $y_{ij}$，则
\begin{equation}
\frac{\partial\mathcal L_i}{\partial s_{ij}}=\frac{p_{ij}-y_{ij}}\tau,
\qquad
\nabla_{u_i}\mathcal L_i=\frac1\tau\left(\sum_jp_{ij}v_j-v_i^+\right).
\end{equation}
梯度提高正例相对分数，并压低有较大当前概率的负例。温度改变概率集中度和梯度尺度，温度应与训练目标共同配置。

例如三个候选的相似度除温度后为 $(\log4,\log2,0)$，第一项为正例，归一化概率与损失为
\begin{equation}
 p=(\frac{4}{7},\frac{2}{7},\frac{1}{7}),\qquad \ell=\log(\frac{7}{4})\approx0.5596.
\end{equation}
对这些缩放后分数的梯度为 $(-\frac{3}{7},\frac{2}{7},\frac{1}{7})$，当前更像正例的第二项获得更大负向学习压力。

\subsection{负例构造及假负例}
\term{批内负例}{In-Batch Negative}{}把其他查询对应文档用作当前查询负例，减少额外编码。\term{难负例}{Hard Negative}{}由词法或现有稠密检索器找出分数高但确实不相关的文档，可以形成更有区分度的学习信号。随机负例通常较容易，过度依赖某一种难例又可能使模型只适应当前检索器的错误。

\term{假负例}{False Negative}{}是真实相关却被标成负例的对象。开放语料中一题往往有多个支持文档，批内文档也可能相互相关；强行推开这些向量会损害召回。可以依据已知相关集合屏蔽冲突、使用多正例损失或对疑似负例复核，但不能把“未标为正”自动理解为“确实不相关”。多正例求和形式还应明确是希望任一正例得高分，还是所有正例都被拉近。

训练与评估应按查询家族、文档版本和必要时间边界隔离。文档原文可以作为开放检索库存在，相关性标注却不能泄漏进查询编码或示例选择。相似文本的近重复治理见\bookxref{ch:data-engineering}。

\section{向量相似度及精确搜索}

\subsection{余弦、内积及欧氏距离}
对非零向量 $u,v$，\term{余弦相似度}{Cosine Similarity}{}为 $\frac{u^{\mathsf T}v}{\|u\|_2\|v\|_2}$。\term{最大内积搜索}{Maximum Inner Product Search}{MIPS}寻找内积最大的向量；\term{欧氏距离}{Euclidean Distance}{}满足
\begin{equation}
\|u-v\|_2^2=\|u\|_2^2+\|v\|_2^2-2u^{\mathsf T}v.
\end{equation}
当文档向量范数固定时，对固定查询最大内积与最小欧氏距离排名一致；若查询和文档均归一化为单位向量，则
\begin{equation}
\|\bar u-\bar v\|_2^2=2-2\cos(u,v).
\label{eq:ch32-metric}
\end{equation}
此时三种排名等价。零向量不能作单位归一化，应隔离或采用明确规则。

未归一化时，范数会改变结果。令 $u=(1,0)$，$v_1=(0.8,0.6)$，$v_2=(2,2)$。余弦分别为 $0.8$ 和 $\frac{1}{\sqrt2}\approx0.7071$，余弦选择 $v_1$；内积分别为 $0.8$ 和 $2$，MIPS 选择 $v_2$。若训练的打分利用范数信息，部署时擅自归一化就改变了目标函数。

\subsection{精确排名基线}
\term{精确最近邻搜索}{Exact Nearest Neighbor Search}{}对所有向量计算指定度量。$N$ 个 $h$ 维向量需要 $O(Nh)$ 评分工作，存储原始单精度向量约 $4Nh$ 字节，尚不含标识和元数据。矩阵批处理可利用高效硬件，但总扫描量仍随 $N$ 增长。

精确搜索给出指定表示与度量下的最近对象，任务相关性通过标注评估。可以将稠密系统误差分成表示误差与索引近似误差：若相关文档在精确前列也不存在，需要改进表示或语料；若精确搜索可以找到但近似索引遗漏，需要调整索引搜索。两者采用不同诊断证据。

\section{近似索引的两条路线}
\optional
\subsection{HNSW 的分层图导航}
\term{近似最近邻搜索}{Approximate Nearest Neighbor Search}{ANN}牺牲部分精确邻居召回，以降低搜索成本。\term{分层可导航小世界图}{Hierarchical Navigable Small World}{HNSW}为向量构建多层邻近图，高层包含较少节点，底层覆盖全部节点\cite{malkov2018hnsw}。查询从稀疏高层快速靠近目标区域，再向下扩展候选；底层保留比贪心单路径更多的搜索状态。

图构建时的邻居预算 $M$、构建搜索宽度与查询搜索宽度影响内存、构建时间和召回。更多边通常提高可导航性，也增加邻接表内存；更宽查询探索更多节点，增加距离计算。启发式邻居选择需要保持方向多样性，若只连接同一局部密簇，图可能难以跨区域导航。

查询工作量由图结构、度量、数据簇和过滤条件决定。增量插入使图结构受构建顺序与随机层级影响，删除或大量墓碑也可能降低导航质量。对高选择性权限过滤，需要考虑合法节点的可达性，不能把过滤后的召回退化归咎于编码器。

\begin{figure}[htbp]
\centering
\begin{tikzpicture}[>=stealth,every node/.style={font=\small},n/.style={circle,draw,minimum size=.45cm,inner sep=1pt}]
\node[anchor=east]at(-.6,1.5){高层};\node[n](a)at(0,1.5){A};\node[n](d)at(6,1.5){D};\draw(a)--(d);
\node[anchor=east]at(-.6,0){底层};
\foreach \name/\x in {aa/0,b/1.5,c/3,dd/4.5,e/6,f/7.5}{\node[n](\name)at(\x,0){};}
\draw(aa)--(b)--(c)--(dd)--(e)--(f);\draw(aa)to[bend left](c);\draw(c)to[bend left](e);
\draw[->,dashed](a)--(aa);\draw[->,dashed](d)--(e);
\node[anchor=west]at(8,0){局部扩展};
\end{tikzpicture}
\caption{HNSW 的概念层级：高层导航与底层候选探索。}
\label{fig:ch32-hnsw}
\end{figure}

\subsection{倒排文件及乘积量化}
\term{倒排文件向量索引}{Inverted File Index}{IVF}先用粗聚类把向量分到若干单元，再将每个单元中的向量标识组织为列表。查询选择有限个最接近的粗中心，只扫描对应列表。它与词法倒排都使用“键到列表”的组织，但键的含义不同：前者是向量单元，后者是词项。

访问更多单元提高候选覆盖，也增加扫描量。若最近向量被分到未访问单元，后续精确重排无法挽回遗漏。粗聚类若在过时或偏移的数据上训练，单元可能失衡，一些查询扫描大量向量而另一些查询漏掉边界邻居。

\term{乘积量化}{Product Quantization}{PQ}把 $h$ 维向量拆成 $m$ 个子向量，每个子空间训练一个含 $k$ 个中心的码本，用 $m$ 个码字索引表示向量\cite{jegou2011pq}。若 $k=256$，每个子码占一字节，向量码为 $m$ 字节，码本及原始精排向量的额外存储另计。

\term{非对称距离计算}{Asymmetric Distance Computation}{ADC}保留查询为高精度，文档用码字中心近似：
\begin{equation}
\|u-\widehat v\|_2^2=\sum_{j=1}^{m}\|u^{(j)}-c_{j,a_j}\|_2^2.
\end{equation}
查询可先建立每个子空间的距离表，再对每条压缩向量查表求和。组合 IVF 与 PQ 时，还可对向量减去粗中心后的残差编码，查询距离必须计入同一个中心，不能把不同单元的残差当成同一原始坐标。

若 $v=\widehat v+e$，反三角不等式给出 $|\|u-v\|_2-\|u-\widehat v\|_2|\le\|e\|_2$。平方距离误差满足
\begin{equation}
\left|\|u-\widehat v\|_2^2-\|u-v\|_2^2\right|
\le2\|u-v\|_2\|e\|_2+\|e\|_2^2.
\end{equation}
小距离间隔的候选更易因近似误差交换名次。IVF 的单元漏检与 PQ 的距离失真是两种误差；扩大扫描单元不能直接消除后者，精排压缩候选也不能恢复前者。

\subsection{索引资源的比较}
HNSW 主要增加图导航结构，IVF 限制访问区域，PQ 压缩表示及距离计算。三种机制可以组合使用。选择时应比较同一表示、同一授权语料和相同精确搜索参照下的 ANN 召回、尾时延、内存、构建成本与更新行为，而非只比较一次查询的平均耗时。

\section{混合检索及学习排序}

\subsection{分数融合的尺度问题}
\term{混合检索}{Hybrid Retrieval}{}结合词法和稠密候选。前者对精确术语、数字标识和罕见字符串有明确匹配依据，后者能够利用训练得到的语义关系；但两者均有失败场景。简单分数和 $\lambda s_{\mathrm{lex}}+(1-\lambda)s_{\mathrm{dense}}$ 只有在尺度处理明确后才有意义。BM25 与余弦的数值范围和分布不同，逐查询最小最大归一化又受异常候选和截断深度影响。

\term{倒数排名融合}{Reciprocal Rank Fusion}{RRF}直接使用排名\cite{cormack2009rrf}：
\begin{equation}
s_{\mathrm{RRF}}(d)=\sum_{j:d\in R_j}\frac{w_j}{c+\operatorname{rank}_j(d)},\qquad c>0.
\end{equation}
未进入某路截断列表的文档在该路贡献为零。$c$ 控制顶部名次差异的影响，$w_j$ 为预先确定的路权重。RRF 不需要跨路原分数可比，却丢失了分数间隔信息；融合收益取决于两路候选的互补程度与相关性质量。

取词法排名 $(A,B,C)$、向量排名 $(B,D,A)$，$c=10,w_j=1$。得到
\begin{align}
s(A)&=\frac{1}{11}+\frac{1}{13}\approx0.16783,\\
s(B)&=\frac{1}{12}+\frac{1}{11}\approx0.17424,\\
s(C)&=\frac{1}{13}\approx0.07692,\qquad s(D)=\frac{1}{12}\approx0.08333.
\end{align}
融合顺序为 $B,A,D,C$。参数 $c$ 可在开发集上选择。相同文档的多个重复块若被当作独立对象，需通过对象身份与文档聚合规则处理冗余。

\subsection{重排及候选集上界}
\term{重排序}{Reranking}{}在有限候选集上使用更精细的评分。\term{交叉编码器}{Cross Encoder}{}联合输入查询与文档，允许细粒度交互，代价是不能像双编码器那样为每篇文档独立预计算最终相关性。\term{学习排序}{Learning to Rank}{LTR}可以组合词法分数、稠密分数、结构和其他合法特征。

例如已标注正负文档对可使用成对损失 $\log(1+\exp[-(s^+-s^-)])$，推动正例高于负例。训练样本来自旧排序器时，需要注意曝光和位置偏差，点击也不等同于完全相关。无论重排器多强，若正确文档不在候选并集中，它都无法将其排到前列；应把召回深度与重排预算共同评价。

\section{权限、版本及增量更新}

\subsection{检索过滤约束}
\term{元数据过滤}{Metadata Filtering}{}约束租户、文档类型、语言、生效时间与访问权限。一个授权查询的检索集合应是当前用户可访问的 $\mathcal D_u$，不是全库评分后让生成器自行忽略秘密。查询时过滤、返回前再次授权检查和缓存隔离应共同保证撤销后的对象不被消费。

只取全库前 $K$ 再删除无权结果，可能返回少于 $K$ 条且漏掉合法相关对象。增加预取深度只能缓解，不提供严格召回保证。ANN 图导航还可能需要经过不符合过滤条件的节点，其标识与内容的内部使用边界应由可信检索服务定义；绝不能向调用方暴露未授权结果。高选择性过滤场景应单独测量授权集合上的召回与时延。

\subsection{索引增量更新}
每条索引对象应绑定稳定文档身份、源版本、片段边界、内容摘要、编码器版本、分析器和索引构建配置。文档编辑可能使后续切分边界整体变化，不能仅按旧片段序号覆盖而保留失效片段。编码器更新则通常要求重编码，并在完整新索引就绪后切换查询编码器与文档空间。

增量处理可记录变更日志，采用幂等写入和不可变版本。\term{删除标记}{Tombstone}{}可先阻止旧条目返回，后续通过物理压缩清理图节点、倒排条目、向量和存储。删除还需传播到片段缓存、检索缓存与派生索引；统计数据与旧快照按各自的保留规则清理。

\term{索引快照}{Index Snapshot}{}将语料版本、编码器、分析器、统计及物理索引绑定为一致消费单位。新版本采用原子可见切换，回滚时恢复匹配的查询配置。权限撤销属于即时约束时，不能因为回滚索引而重新授权已经撤销的内容。

\begin{figure}[htbp]
\centering
\begin{tikzpicture}[>=stealth,every node/.style={font=\small},b/.style={draw,align=center,minimum width=2.35cm,minimum height=.85cm}]
\node[b](doc)at(0,1.2){版本文档};\node[b](chunk)at(3.1,1.2){解析与切分\\身份、边界、权限};\node[b](idx)at(6.4,1.2){词法/向量索引};
\node[b](q)at(0,-1){查询与身份};\node[b](ret)at(3.1,-1){过滤与多路召回};\node[b](rank)at(6.4,-1){融合与重排};\node[b](out)at(9.6,-1){有序授权材料};
\draw[->](doc)--(chunk);\draw[->](chunk)--(idx);\draw[->](idx)--(ret);\draw[->](q)--(ret);\draw[->](ret)--(rank);\draw[->](rank)--(out);
\end{tikzpicture}
\caption{离线语料组织与在线检索链路。生成回答不属于本图的检索职责。}
\label{fig:ch32-pipeline}
\end{figure}

\begin{algorithm}[具有版本边界的混合检索]
\textbf{输入：}查询、已认证身份、一致索引快照、候选深度与排序配置。\textbf{输出：}有序可访问对象及来源。\textbf{状态：}合法过滤条件、各路候选、融合结果。
\begin{enumerate}
\item 解析访问条件与时间范围，选择匹配的分析器和查询编码器；拒绝不兼容版本。
\item 在授权过滤语义下分别执行词法与向量检索，记录真实候选深度及超时状态。
\item 按稳定对象身份合并候选，使用预定分数融合或 RRF，按预算重排并处理同文档冗余。
\item 返回前再次核对当前授权与删除状态，输出实际材料身份、版本、排序分数和诊断状态；资源预算耗尽时按预定规则结束。
\end{enumerate}
\textbf{不变量：}返回内容经过授权，表示与索引版本一致，召回失败或回退不被伪装为正常完整结果。
\end{algorithm}

\section{检索评价指标}

\subsection{覆盖、首条命中及分级排序}
\term{召回率}{Recall}{}衡量相关集合被覆盖的比例：
\begin{equation}
\operatorname{Recall@K}(q)=\frac{|G_q\cap R_q^{(K)}|}{|G_q|},\qquad |G_q|>0.
\end{equation}
\term{命中率}{Hit Rate}{Hit@K}则是 $\mathbf1[G_q\cap R_q^{(K)}\ne\varnothing]$。存在三篇相关文档、只取回一篇时，前者为 $\frac{1}{3}$，后者为 $1$；二者不能混称。没有相关对象的查询分母为零，应单独评价无答案检出或拒绝，而不是擅自记为满分。

\term{平均倒数排名}{Mean Reciprocal Rank}{MRR}关注第一条相关结果。若其位置为 $j_q$，单查询倒数排名为 $\frac{1}{j_q}$，未在规定截断内命中则为零，再对查询平均。它不奖励第一条之后更多相关材料，因此与多证据任务的覆盖目标不同。

\term{折损累计增益}{Discounted Cumulative Gain}{DCG}按等级与位置计算价值，\term{归一化折损累计增益}{Normalized Discounted Cumulative Gain}{nDCG}用理想排序归一化。累积增益评价的基础见 Järvelin 与 Kekäläinen\cite{jarvelin2002dcg}。本章明确采用一种常见变体：
\begin{equation}
\operatorname{DCG@K}=\sum_{j=1}^{K}\frac{2^{r_j}-1}{\log_2(j+1)},\qquad
\operatorname{nDCG@K}=\frac{\operatorname{DCG@K}}{\operatorname{IDCG@K}}.
\end{equation}
理想分母由该查询全部已知相关对象按等级排序取前 $K$ 得到，不能只重排系统已返回对象。若分母为零，应规定无增益查询的独立口径。\footnote{DCG 有线性与指数增益、不同折损位置等约定。比较报告前必须固定公式、相关等级含义和无相关查询处理，不能仅凭相同指标名称认定可比。}

\subsection{同一结果的三种指标}
\begin{example}
\label{q:ch32:example:02}


某查询已知相关对象为 $A,B,C$，等级分别为 $3,2,1$，其余对象等级为零。系统前五名为 $(D,B,E,A,F)$，前四名相关等级为 $(0,2,0,3)$。因此
\begin{equation}
\operatorname{Recall@4}=\frac{2}{3},\quad \operatorname{Hit@4}=1,\quad \operatorname{RR@4}=\frac{1}{2.}
\end{equation}
DCG 为 $\frac{3}{\log_2 3}+\frac{7}{\log_2 5}\approx1.8928+3.0147=4.9075$。理想等级为 $(3,2,1,0)$，IDCG 为 $7+\frac{3}{\log_2 3}+\frac{1}{\log_2 4}\approx9.3928$，所以 nDCG@4 约为 $0.5225$。相关对象 A 虽然已召回，却位于第四位，折损反映其排名损失。

若第二个可回答查询的第一条相关结果在第三位，则两查询 MRR@4 为 $\frac{\frac{1}{2}+\frac{1}{3}}{2}=\frac{5}{12}\approx0.4167$。另有无答案查询时，应单列是否错误返回被当作证据的材料，而不混入上述需要相关集合的分母。
\end{example}

\subsection{ANN 召回及相关性召回分开}
给定精确向量前 $K$ 集合 $E_q^{(K)}$，可定义 ANN 召回为 $\frac{|E_q^{(K)}\cap A_q^{(K)}|}{K}$，用于衡量近似索引对向量排名的保持。它与基于 $G_q$ 的任务召回不同。一个索引可以完美复现不相关的向量排名，也可以因为近似扰动偶然取得更高任务分数；表示质量和索引近似质量应分别评价。

标注不完整时，未标注文档不一定不相关，召回分母与 nDCG 理想排序均只反映当前已知集合。应保存标注来源，进行多路候选池复核，并按精确实体、语义改写、跨语言、长文档、时效、权限过滤和无答案查询等难例切片报告。宏平均让每个查询等权，按相关对象数加权则回答不同问题。

\begin{note}[检索故障的定位顺序]
先确认解析与权限范围，再检查相关材料是否进入精确候选，随后检查近似索引与融合是否遗漏，最后检查重排是否把有效材料移出前列。召回、排名与下游答案质量分别评价候选覆盖、排序与最终使用效果。
\end{note}


% BEGIN GENERATED INTERVIEWS

% END GENERATED INTERVIEWS
\chapterreferences
\end{document}
